\IfFileExists{stacks-project.cls}{%
\documentclass{stacks-project}
}{%
\documentclass{amsart}
}

% For dealing with references we use the comment environment
\usepackage{verbatim}
\newenvironment{reference}{\comment}{\endcomment}
%\newenvironment{reference}{}{}
\newenvironment{slogan}{\comment}{\endcomment}
\newenvironment{history}{\comment}{\endcomment}

% For commutative diagrams we use Xy-pic
\usepackage[all]{xy}

% We use 2cell for 2-commutative diagrams.
\xyoption{2cell}
\UseAllTwocells

% We use multicol for the list of chapters between chapters
\usepackage{multicol}

% This is generally recommended for better output
\usepackage{lmodern}
\usepackage[T1]{fontenc}

% For cross-file-references
\usepackage{xr-hyper}

% Package for hypertext links:
\usepackage{hyperref}

% For any local file, say "hello.tex" you want to link to please
% use \externaldocument[hello-]{hello}
\externaldocument[introduction-]{introduction}
\externaldocument[conventions-]{conventions}
\externaldocument[sets-]{sets}
\externaldocument[categories-]{categories}
\externaldocument[topology-]{topology}
\externaldocument[sheaves-]{sheaves}
\externaldocument[sites-]{sites}
\externaldocument[stacks-]{stacks}
\externaldocument[fields-]{fields}
\externaldocument[algebra-]{algebra}
\externaldocument[brauer-]{brauer}
\externaldocument[homology-]{homology}
\externaldocument[derived-]{derived}
\externaldocument[simplicial-]{simplicial}
\externaldocument[more-algebra-]{more-algebra}
\externaldocument[smoothing-]{smoothing}
\externaldocument[modules-]{modules}
\externaldocument[sites-modules-]{sites-modules}
\externaldocument[injectives-]{injectives}
\externaldocument[cohomology-]{cohomology}
\externaldocument[sites-cohomology-]{sites-cohomology}
\externaldocument[dga-]{dga}
\externaldocument[dpa-]{dpa}
\externaldocument[sdga-]{sdga}
\externaldocument[hypercovering-]{hypercovering}
\externaldocument[schemes-]{schemes}
\externaldocument[constructions-]{constructions}
\externaldocument[properties-]{properties}
\externaldocument[morphisms-]{morphisms}
\externaldocument[coherent-]{coherent}
\externaldocument[divisors-]{divisors}
\externaldocument[limits-]{limits}
\externaldocument[varieties-]{varieties}
\externaldocument[topologies-]{topologies}
\externaldocument[descent-]{descent}
\externaldocument[perfect-]{perfect}
\externaldocument[more-morphisms-]{more-morphisms}
\externaldocument[flat-]{flat}
\externaldocument[groupoids-]{groupoids}
\externaldocument[more-groupoids-]{more-groupoids}
\externaldocument[etale-]{etale}
\externaldocument[chow-]{chow}
\externaldocument[intersection-]{intersection}
\externaldocument[pic-]{pic}
\externaldocument[weil-]{weil}
\externaldocument[adequate-]{adequate}
\externaldocument[dualizing-]{dualizing}
\externaldocument[duality-]{duality}
\externaldocument[discriminant-]{discriminant}
\externaldocument[derham-]{derham}
\externaldocument[local-cohomology-]{local-cohomology}
\externaldocument[algebraization-]{algebraization}
\externaldocument[curves-]{curves}
\externaldocument[resolve-]{resolve}
\externaldocument[models-]{models}
\externaldocument[functors-]{functors}
\externaldocument[equiv-]{equiv}
\externaldocument[pione-]{pione}
\externaldocument[etale-cohomology-]{etale-cohomology}
\externaldocument[proetale-]{proetale}
\externaldocument[relative-cycles-]{relative-cycles}
\externaldocument[more-etale-]{more-etale}
\externaldocument[trace-]{trace}
\externaldocument[crystalline-]{crystalline}
\externaldocument[spaces-]{spaces}
\externaldocument[spaces-properties-]{spaces-properties}
\externaldocument[spaces-morphisms-]{spaces-morphisms}
\externaldocument[decent-spaces-]{decent-spaces}
\externaldocument[spaces-cohomology-]{spaces-cohomology}
\externaldocument[spaces-limits-]{spaces-limits}
\externaldocument[spaces-divisors-]{spaces-divisors}
\externaldocument[spaces-over-fields-]{spaces-over-fields}
\externaldocument[spaces-topologies-]{spaces-topologies}
\externaldocument[spaces-descent-]{spaces-descent}
\externaldocument[spaces-perfect-]{spaces-perfect}
\externaldocument[spaces-more-morphisms-]{spaces-more-morphisms}
\externaldocument[spaces-flat-]{spaces-flat}
\externaldocument[spaces-groupoids-]{spaces-groupoids}
\externaldocument[spaces-more-groupoids-]{spaces-more-groupoids}
\externaldocument[bootstrap-]{bootstrap}
\externaldocument[spaces-pushouts-]{spaces-pushouts}
\externaldocument[spaces-chow-]{spaces-chow}
\externaldocument[groupoids-quotients-]{groupoids-quotients}
\externaldocument[spaces-more-cohomology-]{spaces-more-cohomology}
\externaldocument[spaces-simplicial-]{spaces-simplicial}
\externaldocument[spaces-duality-]{spaces-duality}
\externaldocument[formal-spaces-]{formal-spaces}
\externaldocument[restricted-]{restricted}
\externaldocument[spaces-resolve-]{spaces-resolve}
\externaldocument[formal-defos-]{formal-defos}
\externaldocument[defos-]{defos}
\externaldocument[cotangent-]{cotangent}
\externaldocument[examples-defos-]{examples-defos}
\externaldocument[algebraic-]{algebraic}
\externaldocument[examples-stacks-]{examples-stacks}
\externaldocument[stacks-sheaves-]{stacks-sheaves}
\externaldocument[criteria-]{criteria}
\externaldocument[artin-]{artin}
\externaldocument[quot-]{quot}
\externaldocument[stacks-properties-]{stacks-properties}
\externaldocument[stacks-morphisms-]{stacks-morphisms}
\externaldocument[stacks-limits-]{stacks-limits}
\externaldocument[stacks-cohomology-]{stacks-cohomology}
\externaldocument[stacks-perfect-]{stacks-perfect}
\externaldocument[stacks-introduction-]{stacks-introduction}
\externaldocument[stacks-more-morphisms-]{stacks-more-morphisms}
\externaldocument[stacks-geometry-]{stacks-geometry}
\externaldocument[moduli-]{moduli}
\externaldocument[moduli-curves-]{moduli-curves}
\externaldocument[examples-]{examples}
\externaldocument[exercises-]{exercises}
\externaldocument[guide-]{guide}
\externaldocument[desirables-]{desirables}
\externaldocument[coding-]{coding}
\externaldocument[obsolete-]{obsolete}
\externaldocument[fdl-]{fdl}
\externaldocument[index-]{index}

% Theorem environments.
%
\theoremstyle{plain}
\newtheorem{theorem}[subsection]{Theorem}
\newtheorem{proposition}[subsection]{Proposition}
\newtheorem{lemma}[subsection]{Lemma}

\theoremstyle{definition}
\newtheorem{definition}[subsection]{Definition}
\newtheorem{example}[subsection]{Example}
\newtheorem{exercise}[subsection]{Exercise}
\newtheorem{situation}[subsection]{Situation}

\theoremstyle{remark}
\newtheorem{remark}[subsection]{Remark}
\newtheorem{remarks}[subsection]{Remarks}

\numberwithin{equation}{subsection}

% Macros
%
\def\lim{\mathop{\mathrm{lim}}\nolimits}
\def\colim{\mathop{\mathrm{colim}}\nolimits}
\def\Spec{\mathop{\mathrm{Spec}}}
\def\Hom{\mathop{\mathrm{Hom}}\nolimits}
\def\Ext{\mathop{\mathrm{Ext}}\nolimits}
\def\SheafHom{\mathop{\mathcal{H}\!\mathit{om}}\nolimits}
\def\SheafExt{\mathop{\mathcal{E}\!\mathit{xt}}\nolimits}
\def\Sch{\mathit{Sch}}
\def\Mor{\mathop{\mathrm{Mor}}\nolimits}
\def\Ob{\mathop{\mathrm{Ob}}\nolimits}
\def\Sh{\mathop{\mathit{Sh}}\nolimits}
\def\NL{\mathop{N\!L}\nolimits}
\def\CH{\mathop{\mathrm{CH}}\nolimits}
\def\proetale{{pro\text{-}\acute{e}tale}}
\def\etale{{\acute{e}tale}}
\def\QCoh{\mathit{QCoh}}
\def\Ker{\mathop{\mathrm{Ker}}}
\def\Im{\mathop{\mathrm{Im}}}
\def\Coker{\mathop{\mathrm{Coker}}}
\def\Coim{\mathop{\mathrm{Coim}}}

% Boxtimes
%
\DeclareMathSymbol{\boxtimes}{\mathbin}{AMSa}{"02}

%
% Macros for moduli stacks/spaces
%
\def\QCohstack{\mathcal{QC}\!\mathit{oh}}
\def\Cohstack{\mathcal{C}\!\mathit{oh}}
\def\Spacesstack{\mathcal{S}\!\mathit{paces}}
\def\Quotfunctor{\mathrm{Quot}}
\def\Hilbfunctor{\mathrm{Hilb}}
\def\Curvesstack{\mathcal{C}\!\mathit{urves}}
\def\Polarizedstack{\mathcal{P}\!\mathit{olarized}}
\def\Complexesstack{\mathcal{C}\!\mathit{omplexes}}
% \Pic is the operator that assigns to X its picard group, usage \Pic(X)
% \Picardstack_{X/B} denotes the Picard stack of X over B
% \Picardfunctor_{X/B} denotes the Picard functor of X over B
\def\Pic{\mathop{\mathrm{Pic}}\nolimits}
\def\Picardstack{\mathcal{P}\!\mathit{ic}}
\def\Picardfunctor{\mathrm{Pic}}
\def\Deformationcategory{\mathcal{D}\!\mathit{ef}}
\usepackage{cleveref}
\crefname{originalreference}{original reference}{original references}
\Crefname{originalreference}{Original reference}{Original references}
\makeatletter
\newlabel{algebra-lemma-characterize-finite-presentation}{{00QO}{source}{}{sourcepointer-1}{}}
\newlabel{algebra-lemma-characterize-finite-presentation@cref}{{[originalreference][0][]00QO}{[1][0][]source}{}{}{}}
\makeatother
\crefname{originaldefinition}{original definition}{original definitions}
\Crefname{originaldefinition}{Original definition}{Original definitions}
\makeatletter
\newlabel{definition-locally-finite-presentation}{{049J}{source}{}{sourcepointer-2}{}}
\newlabel{definition-locally-finite-presentation@cref}{{[originaldefinition][0][]049J}{[1][0][]source}{}{}{}}
\makeatother
\crefname{originallemma}{original lemma}{original lemmas}
\Crefname{originallemma}{Original lemma}{Original lemmas}
\makeatletter
\newlabel{lemma-base-change-locally-finite-presentation}{{049M}{source}{}{sourcepointer-3}{}}
\newlabel{lemma-base-change-locally-finite-presentation@cref}{{[originallemma][0][]049M}{[1][0][]source}{}{}{}}
\makeatother
\crefname{originallemma}{original lemma}{original lemmas}
\Crefname{originallemma}{Original lemma}{Original lemmas}
\makeatletter
\newlabel{lemma-composition-locally-finite-presentation}{{049L}{source}{}{sourcepointer-4}{}}
\newlabel{lemma-composition-locally-finite-presentation@cref}{{[originallemma][0][]049L}{[1][0][]source}{}{}{}}
\makeatother
\crefname{originallemma}{original lemma}{original lemmas}
\Crefname{originallemma}{Original lemma}{Original lemmas}
\makeatletter
\newlabel{lemma-sheafify-finite-presentation}{{049O}{source}{}{sourcepointer-5}{}}
\newlabel{lemma-sheafify-finite-presentation@cref}{{[originallemma][0][]049O}{[1][0][]source}{}{}{}}
\makeatother
\crefname{originalreference}{original reference}{original references}
\Crefname{originalreference}{Original reference}{Original references}
\makeatletter
\newlabel{limits-proposition-characterize-locally-finite-presentation}{{04AK}{source}{}{sourcepointer-6}{}}
\newlabel{limits-proposition-characterize-locally-finite-presentation@cref}{{[originalreference][0][]04AK}{[1][0][]source}{}{}{}}
\makeatother
\crefname{originalreference}{original reference}{original references}
\Crefname{originalreference}{Original reference}{Original references}
\makeatletter
\newlabel{spaces-definition-presentation}{{0263}{source}{}{sourcepointer-7}{}}
\newlabel{spaces-definition-presentation@cref}{{[originalreference][0][]0263}{[1][0][]source}{}{}{}}
\makeatother
\crefname{originalreference}{original reference}{original references}
\Crefname{originalreference}{Original reference}{Original references}
\makeatletter
\newlabel{spaces-morphisms-definition-locally-finite-presentation}{{03XP}{source}{}{sourcepointer-8}{}}
\newlabel{spaces-morphisms-definition-locally-finite-presentation@cref}{{[originalreference][0][]03XP}{[1][0][]source}{}{}{}}
\makeatother
\crefname{originalreference}{original reference}{original references}
\Crefname{originalreference}{Original reference}{Original references}
\makeatletter
\newlabel{spaces-morphisms-lemma-etale-locally-finite-presentation}{{0468}{source}{}{sourcepointer-9}{}}
\newlabel{spaces-morphisms-lemma-etale-locally-finite-presentation@cref}{{[originalreference][0][]0468}{[1][0][]source}{}{}{}}
\makeatother
\title{Locally finite presentation equals limit preservation for algebraic spaces}
\author{The Stacks Project Authors}
\date{stacks-project:spaces-limits.tex:local-snapshot-20261012; GNU FDL 1.2 or later}
\begin{document}
\maketitle
\noindent\textit{Editorial note.} Counted span is the author's own proof environment, source lines 708-781, following the statement at 693-706.\par\medskip
\noindent\textit{Editorial note.} No omissions flagged; the set-theoretic remark is a complete argument, not an aside.\par\medskip
\noindent\textit{Allowed standard prerequisites (editorial).}\begin{itemize}
\item Presentation of an algebraic space. Spaces, Definition spaces-definition-presentation (cross-chapter, cited at \textasciitilde{}716-717). Chooses X = U/R.
\item Scheme-level limit preservation characterization. Limits, Proposition limits-proposition-characterize-locally-finite-presentation (cross-chapter, cited at \textasciitilde{}723-724, \textasciitilde{}755-756, \textasciitilde{}763-773). U(T) = colim U(T\_i), R(T) = colim R(T\_i); and the converse for schemes.
\item Sheafification preserves limit preservation. lemma-sheafify-finite-presentation (same chapter, cited at \textasciitilde{}737-738). Passes from the presheaf W |-> U(W)/R(W) to X = U/R.
\item Limit preservation under base change and composition. lemma-base-change-locally-finite-presentation and lemma-composition-locally-finite-presentation (same chapter, cited at \textasciitilde{}744-745 and \textasciitilde{}751-753). Propagates limit preservation along the etale covers in direction (2).
\item Etale maps are locally of finite presentation. Morphisms of Spaces, Lemma spaces-morphisms-lemma-etale-locally-finite-presentation (cross-chapter, cited at \textasciitilde{}747-749). U -> V x\_Y X is locally of finite presentation.
\end{itemize}\medskip
\ifdefined\hypertarget\hypertarget{sourcepointer-1}{}\fi
\noindent\textit{Source reference.} Original reference 00QO: focused-seeds/geometry/S9-spaces-limits-characterize-locally-finite-presentation/ref-tags.txt, source lines 5--5. Navigation pointer only: the counted proof cites this statement by \textbackslash{}ref as a standing input; the deduction the counted proof has to make is carried out inside the counted span, and no part of that deduction is delegated to this pointer.\par\medskip
\ifdefined\hypertarget\hypertarget{sourcepointer-2}{}\fi
\noindent\textit{Source reference.} Original definition 049J: sources/stacks-spaces-limits/pinned-original-source.tex, source lines 287--287. Navigation pointer only: the counted proof cites this statement by \textbackslash{}ref as a standing input; the deduction the counted proof has to make is carried out inside the counted span, and no part of that deduction is delegated to this pointer.\par\medskip
\ifdefined\hypertarget\hypertarget{sourcepointer-3}{}\fi
\noindent\textit{Source reference.} Original lemma 049M: sources/stacks-spaces-limits/pinned-original-source.tex, source lines 451--451. Navigation pointer only: the counted proof cites this statement by \textbackslash{}ref as a standing input; the deduction the counted proof has to make is carried out inside the counted span, and no part of that deduction is delegated to this pointer.\par\medskip
\ifdefined\hypertarget\hypertarget{sourcepointer-4}{}\fi
\noindent\textit{Source reference.} Original lemma 049L: sources/stacks-spaces-limits/pinned-original-source.tex, source lines 399--399. Navigation pointer only: the counted proof cites this statement by \textbackslash{}ref as a standing input; the deduction the counted proof has to make is carried out inside the counted span, and no part of that deduction is delegated to this pointer.\par\medskip
\ifdefined\hypertarget\hypertarget{sourcepointer-5}{}\fi
\noindent\textit{Source reference.} Original lemma 049O: sources/stacks-spaces-limits/pinned-original-source.tex, source lines 506--506. Navigation pointer only: the counted proof cites this statement by \textbackslash{}ref as a standing input; the deduction the counted proof has to make is carried out inside the counted span, and no part of that deduction is delegated to this pointer.\par\medskip
\ifdefined\hypertarget\hypertarget{sourcepointer-6}{}\fi
\noindent\textit{Source reference.} Original reference 04AK: focused-seeds/geometry/S9-spaces-limits-characterize-locally-finite-presentation/ref-tags.txt, source lines 6--6. Navigation pointer only: the counted proof cites this statement by \textbackslash{}ref as a standing input; the deduction the counted proof has to make is carried out inside the counted span, and no part of that deduction is delegated to this pointer.\par\medskip
\ifdefined\hypertarget\hypertarget{sourcepointer-7}{}\fi
\noindent\textit{Source reference.} Original reference 0263: focused-seeds/geometry/S9-spaces-limits-characterize-locally-finite-presentation/ref-tags.txt, source lines 7--7. Navigation pointer only: the counted proof cites this statement by \textbackslash{}ref as a standing input; the deduction the counted proof has to make is carried out inside the counted span, and no part of that deduction is delegated to this pointer.\par\medskip
\ifdefined\hypertarget\hypertarget{sourcepointer-8}{}\fi
\noindent\textit{Source reference.} Original reference 03XP: focused-seeds/geometry/S9-spaces-limits-characterize-locally-finite-presentation/ref-tags.txt, source lines 8--8. Navigation pointer only: the counted proof cites this statement by \textbackslash{}ref as a standing input; the deduction the counted proof has to make is carried out inside the counted span, and no part of that deduction is delegated to this pointer.\par\medskip
\ifdefined\hypertarget\hypertarget{sourcepointer-9}{}\fi
\noindent\textit{Source reference.} Original reference 0468: focused-seeds/geometry/S9-spaces-limits-characterize-locally-finite-presentation/ref-tags.txt, source lines 9--9. Navigation pointer only: the counted proof cites this statement by \textbackslash{}ref as a standing input; the deduction the counted proof has to make is carried out inside the counted span, and no part of that deduction is delegated to this pointer.\par\medskip
\begin{proposition}
\label{proposition-characterize-locally-finite-presentation}
Let $S$ be a scheme. Let $f : X \to Y$ be a morphism of algebraic
spaces over $S$. The following are equivalent:
\begin{enumerate}
\item The morphism $f$ is a morphism of algebraic spaces which is
locally of finite presentation, see
Morphisms of Spaces,
Definition \ref{spaces-morphisms-definition-locally-finite-presentation}.
\item The morphism $f : X \to Y$ is limit preserving as
a transformation of functors, see
Definition \ref{definition-locally-finite-presentation}.
\end{enumerate}
\end{proposition}

\begin{proof}
Assume (1). Let $T$ be a scheme and let $y \in Y(T)$. We have to show that
$T \times_Y X$ is limit preserving over $T$ in the sense of
Definition \ref{definition-locally-finite-presentation}.
Hence we are reduced to proving that if $X$ is an algebraic space which
is locally of finite presentation over $S$ as an algebraic space, then it
is limit preserving as a functor
$X : (\Sch/S)_{fppf}^{opp} \to \textit{Sets}$.
To see this choose a presentation $X = U/R$, see
Spaces, Definition \ref{spaces-definition-presentation}.
It follows from
Morphisms of Spaces,
Definition \ref{spaces-morphisms-definition-locally-finite-presentation}
that both $U$ and $R$ are schemes which are locally of finite presentation
over $S$. Hence by
Limits, Proposition
\ref{limits-proposition-characterize-locally-finite-presentation}
we have
$$
U(T) = \colim U(T_i), \quad
R(T) = \colim R(T_i)
$$
whenever $T = \lim_i T_i$ in $(\Sch/S)_{fppf}$. It follows
that the presheaf
$$
(\Sch/S)_{fppf}^{opp} \longrightarrow \textit{Sets}, \quad
W \longmapsto U(W)/R(W)
$$
is limit preserving. Hence by
Lemma \ref{lemma-sheafify-finite-presentation}
its sheafification $X = U/R$ is limit preserving too.

\medskip\noindent
Assume (2). Choose a scheme $V$ and a surjective \'etale morphism
$V \to Y$. Next, choose a scheme $U$ and a surjective \'etale morphism
$U \to V \times_Y X$. By
Lemma \ref{lemma-base-change-locally-finite-presentation}
the transformation of functors $V \times_Y X \to V$ is limit preserving. By
Morphisms of Spaces,
Lemma \ref{spaces-morphisms-lemma-etale-locally-finite-presentation}
the morphism of algebraic spaces $U \to V \times_Y X$ is locally
of finite presentation, hence limit preserving as
a transformation of functors by the first part of the proof. By
Lemma \ref{lemma-composition-locally-finite-presentation}
the composition $U \to V \times_Y X \to V$ is limit preserving
as a transformation of functors. Hence
the morphism of schemes $U \to V$ is locally of finite presentation by
Limits, Proposition
\ref{limits-proposition-characterize-locally-finite-presentation}
(modulo a set theoretic remark, see last paragraph of the proof).
This means, by definition, that (1) holds.

\medskip\noindent
Set theoretic remark. Let $U \to V$ be a morphism of
$(\Sch/S)_{fppf}$. In the statement of
Limits, Proposition
\ref{limits-proposition-characterize-locally-finite-presentation}
we characterize $U \to V$ as being locally of finite presentation
if for {\it all} directed inverse systems $(T_i, f_{ii'})$ of affine schemes
over $V$ we have $U(T) = \colim V(T_i)$, but in the current setting
we may only consider affine schemes $T_i$ over $V$ which are (isomorphic to)
an object of $(\Sch/S)_{fppf}$. So we have to make sure that there
are enough affines in $(\Sch/S)_{fppf}$ to make the proof work.
Inspecting the proof of (2) $\Rightarrow$ (1) of
Limits, Proposition
\ref{limits-proposition-characterize-locally-finite-presentation}
we see that the question reduces to the case that $U$ and $V$ are affine.
Say $U = \Spec(A)$ and $V = \Spec(B)$. By construction
of $(\Sch/S)_{fppf}$ the spectrum of any ring of cardinality
$\leq |B|$ is isomorphic to an object of $(\Sch/S)_{fppf}$.
Hence it suffices to observe that in the "only if" part of the proof of
Algebra, Lemma \ref{algebra-lemma-characterize-finite-presentation}
only $A$-algebras of cardinality $\leq |B|$ are used.
\end{proof}

\end{document}
